\section{Computer-assisted certificates and reproducibility}\label{d:sec:reproducibility}

Two finite computer-assisted proofs enter the construction: the
spectral certificate of the period-20 Yee background (Part~III) and the
corner certificate of the continuum background (Part~II). Each checks
finitely many exact inequalities; the analytic arguments of the
respective part extend them to the infinite-dimensional operators. The
logical compiler of Part~I uses no certificate. Its worked example in
Figure~\ref{d:fig:overview} is an exact finite computation.

\subsection{The Yee spectral certificate}\label{si:sec:yee-certificate}

The spectral construction uses exact finite inequalities for the specified
periodic Maxwell operator. Its certificate data comprise 604 bulk-gap
matrices, 2056 global-comparator matrices, local point and neighborhood
data, and a reciprocal-space cover. The matrix and local estimates enter
the spectral assembly in Section~\ref{d:app:certificate}; the resulting
band edge is the premise for the analytic defect and interaction proofs.

The checker reconstructs the incidence and material operators, symmetry
maps, interval inequalities and geometric coverage from the certificate
inputs. Integer and rational operations are exact. Interval computations
use outward bounds with their rounding errors included. Proposed
enclosures are supplied as data arrays, and the required matrix
inequalities are tested against the reconstructed operator.

In summary, the Yee certificate checks:
\begin{itemize}
\item the exact integer material table~\eqref{d:eq:physical-exact-material},
its primitive translations and its scalar inversion symmetry;
\item the eigenvalue enclosures and inertia counts at the four
inequivalent $\mathrm{L}$ points, the center ordering that selects $\mathrm{L}_{++-}$,
and the two-mode cluster data;
\item the local matrix inequalities~\eqref{d:eq:appendix-local-matrix-bounds}
and the uniform inertia count on each radius-$1/20$ neighborhood;
\item the bulk-gap and global-comparator interval matrices and the exact
polytope cover of the irreducible wedge, which give the global
margin~\eqref{d:eq:appendix-global-margin}.
\end{itemize}
All groups are needed for the global spectral conclusion.

\paragraph{Status.} The checker has been run on the complete certificate data,
and all finite inequalities listed above pass. The final step of the argument,
i.e., a complete and independently verified binding of the checked point,
center, derivative and global matrix data to the exact Yee operator, including
the argument in the punctured neighborhoods of the selected band minimum, is
being completed as part of an ongoing formal verification. Until this step is
complete, Theorem~\ref{d:thm:main} and the Yee transfer
Theorem~\ref{d:thm:algorithm-transfer} are conditional on it. The continuum
results of Part~II do not depend on the Yee certificate.

\subsection{The continuum corner certificate}\label{c:app:certificate}

The continuum checker is a standard-library program in exact rational
arithmetic. It represents $s+t\sqrt{19}$ by the exact rational pair
$(s,t)$. Multiplication uses $\sqrt{19}^2=19$, and inversion uses
$(s-t\sqrt{19})/(s^2-19t^2)$ after checking a nonzero denominator.
It reconstructs the 24-component corner projection $P_\varsigma$ and the
actual neighboring-wavevector coupling $-2^{j-2}P_\varsigma P_{\varsigma'}$.
For each of the eight reflection parities it verifies its
intertwining with~\eqref{c:eq:parityblock}, projection invariance,
the Gram normalization, pairwise parity orthogonality, and the
trace and trace-square identities
for the two transverse roots. The velocity Gram identities are
checked on the reconstructed corner vector, not assumed from a
list of eigenvalues. It then checks the three exact Hessian
coefficients~\eqref{c:eq:explicitA0} by inversion in
$\Q(\sqrt{19})$.

For the order and placement comparisons, rational interval
arithmetic encloses all sixteen roots in~\eqref{c:eq:clusterroots},
$\pi$, and the necessary square roots. Square-root endpoints are
obtained by integer square root after scaling by $2^{440}$;
rounding the lower endpoint down and the upper endpoint up
rigorously encloses the true value. Machin's formula uses 128 terms
of each alternating arctangent series and the next-term error.
Logarithms use
\[
 \log s=2\sum_{j=0}^{127}\frac{w^{2j+1}}{2j+1}+R,
 \qquad w=\frac{s-1}{s+1},\qquad
 |R|\le\frac{2|w|^{257}}{257(1-w^2)}.
\]
This formula follows by integrating the geometric series and is
used only for $s=2,3$. Every interval operation is exact rational
endpoint arithmetic. Induction on the operations proves enclosure;
no machine floating-point value enters the calculation.

The accepted inequalities verify the cluster gap $>1/20$, the
finite-contour perturbation inequalities in Lemma~\ref{c:lem:edge},
the Hessian and quartic bounds used there, the fixed integer
placement~\eqref{c:eq:Bmetric} including the $2^{-60}$ Hessian error,
and the strict rate comparisons
$2(.99)^2/(7/4)>\log3$, which with $\sqrt3<7/4$ gives
$2(.99)^2/\sqrt3>\log3$, and
$\log8-2(1.01)^2>1/50$. The verifier contains all its finite input
data. It also checks the small-integer logarithmic inequality for the
fixed majorant ledger of Section~\ref{c:app:ledger}. It uses no
downloaded data, mesh eigensolver or special-function package.

These finite algebraic inequalities enter the continuum estimates at the
places specified above. The resolvent, auxiliary-mode, interaction,
leakage and boundary arguments control their infinite-dimensional
extensions. Sections~\ref{c:sec:parameters} and~\ref{c:sec:cost} give the
terminating parameter choices and account for the complete encoding cost.

\subsection{Rerunning the checks}

The accompanying research materials contain both checkers, their
complete input data, and the data behind the figures. The Yee checker
has a single entry point that runs the matrix, local and coverage
checks; an input manifest identifies the coefficient and enclosure data
used by each check. The continuum checker requires only a standard
Python interpreter and its standard library. It must be run with
assertions enabled, and it refuses to run otherwise. On success it
prints a final confirmation line stating that all exact rational checks
passed and writes a receipt recording the rational intervals and exact
algebraic identities. The receipt can be regenerated and is not a
substitute for running the checks.

The exact logical example in Figure~\ref{d:fig:overview} includes its sparse
matrices, paths, residuals and decoded coordinates.
Figures~\ref{d:fig:FIG-2} and~\ref{d:fig:FIG-3} use the material and spectral
certificate data.
The geometric and defect drawings are schematic.
